\section{Interface and learning objective}

Let $q$ be a Japanese question and let $P(q)$ be its official PLaMo translation-request prompt, with Japanese as the source language and English as the requested output language.
The prompt ends at the start of the translation output; it contains neither a reference translation nor an answer.
At source layer $\ell$, a forward pass gives states
$H^{(\ell)} \in \mathbb{R}^{T\times d_s}$ for the $T$ valid prompt tokens.
The selected source layers are 4, 8 and 16, with layer 8 used for the main configuration.
PLaMo's causal attention means that the final valid state is computed after the complete question, while earlier states have access only to preceding input.

Ordered contiguous bins compress these states to a prefix of length $K$.
For bin $j=0,\ldots,K-1$, define
\begin{equation}
 a_j=\left\lfloor\frac{jT}{K}\right\rfloor,
 \qquad
 b_j=\left\lfloor\frac{(j+1)T}{K}\right\rfloor,
 \qquad
 \bar h_j=\frac{1}{b_j-a_j}\sum_{t=a_j}^{b_j-1}H^{(\ell)}_t.
 \label{eq:pooling}
\end{equation}
Equation~\ref{eq:pooling} applies when $T\geq K$, as in the experimental prompts.
The implementation repeats the corresponding valid state for an empty bin when $T<K$.
Padding is removed before binning; translation-format tokens remain included.
The last bin includes the final valid state.
The main prefix has $K=16$, with $K=4$ as an ablation.

A shared affine map produces target-space embeddings:
\begin{equation}
 z_j=W\bar h_j+b,
 \qquad W\in\mathbb{R}^{d_t\times d_s},
 \quad b\in\mathbb{R}^{d_t}.
 \label{eq:bridge}
\end{equation}
Both models have hidden size 1,536 in the selected configurations, so this bridge has 2,360,832 trainable parameters.
Equal dimensions do not make their representation spaces interchangeable; $W$ and $b$ are learned.
The projected prefix replaces a temporary marker inside the target chat template's user message.
The rest of the prompt asks for the final numeric answer, and the normal assistant-generation boundary follows it.
The marker itself is not embedded by the target model.
The target receives no Japanese question tokens on the bridge path.

For answer tokens $y_1,\ldots,y_M$, including the termination token, training minimizes
\begin{equation}
 \mathcal{L}(W,b)
 =-\frac{1}{M}\sum_{m=1}^{M}
 \log p_{\theta}\!\left(y_m\mid U_1,Z,U_2,y_{<m}\right),
 \label{eq:loss}
\end{equation}
where $U_1$ and $U_2$ are the fixed template and instruction segments, $Z=(z_0,\ldots,z_{K-1})$, and $\theta$ denotes the frozen target parameters.
Prompt and padding positions are excluded from the loss.
The target forward pass remains differentiable so that gradients reach the bridge through the frozen target.
Source states are cached without gradients.
At evaluation, greedy decoding starts at the same assistant boundary without answer tokens or teacher forcing.

The supervised textual baseline replaces $Z$ with $K$ learned embeddings shared across questions and keeps the Japanese question in ordinary target text.
At $K=16$ it has 24,576 trainable parameters.
It therefore matches the training examples but neither the parameter count nor the information path of the bridge.
